\section{Conclusion}
\label{sec:conclusion}
We study apparent prompt authorship through controlled text contrasts, separate-session source judgments, and local activation interventions. Exact-content formatting changes Gemini's elicited authorship classifications substantially, but none of six primary behavioral comparisons is significant after correction. Broader paraphrases fail semantic controls at scale; local activation separability is largely matched by surface features; and fixed-text edits do not establish a style-specific behavioral mechanism. The general human-versus-LLM authorship hypothesis remains open.

Future studies should use human-audited proposition-preserving pairs or richer fixed-content factorials, include borderline-risk requests, and power analyses from observed discordance. The Qwen unconditional-alignment signal requires replication before interpretation. If a natural effect replicates, necessity and matched-coordinate rescue should be tested across layers and tokens with more random and nuisance controls and independently validated source representations.
